\section{Khái niệm đặt vào trong mô hình}
\label{sec:ch15-khai-niem-trong-mo-hinh}

\index{mô hình nút thắt khái niệm!kiến trúc|(}%
Phần mở đầu chương đã nói lối thoát của Phần~V nằm ở chỗ nào. Mục này dựng lại
kiến trúc thực hiện lối thoát ấy, vì mọi thứ còn lại của chương, kể cả chỗ hỏng,
đều đọc từ ba chi tiết của kiến trúc: cái gì được trưng ra, cái gì được huấn
luyện, và cái gì đi qua đâu.

Một \tn{mô hình nút thắt khái niệm}{concept bottleneck model}, từ đây viết tắt
là CBM, nhận đầu vào $x$ và trả nhãn dự đoán $\hat y$ qua đúng hai chặng. Chặng thứ nhất là bộ mã hóa khái niệm, ánh xạ $x$ thành một vector
các giá trị kích hoạt. Chặng thứ hai là đầu phân loại, nhận vector ấy và trả
$\hat y$. Cái làm chặng thứ nhất khác một tầng ẩn thông thường là nó có nhãn
riêng: bộ dữ liệu, ngoài nhãn nhiệm vụ $y$, còn kèm một tập khái niệm $\mathcal{C}$
do con người gán, và với mỗi điểm dữ liệu thì mỗi khái niệm $k$ trong $\mathcal{C}$
có một giá trị thật $c_k$. Sách viết $\hat c_k$ cho giá trị mà bộ mã hóa dự đoán
cho khái niệm ấy. Trong ba bài của chương, $c_k$ luôn nhị phân, nghĩa là khái
niệm ấy có mặt hay vắng mặt ở điểm dữ liệu đang xét.

\index{mô hình nút thắt khái niệm!ví dụ y sinh}%
Ví dụ mà bài của mục~\ref{sec:ch15-do-ro-ri} dùng làm cảnh nền là chẩn đoán ung
thư: $x$ là ảnh sinh thiết, các khái niệm là những thứ bác sĩ vẫn đọc ra từ ảnh
ấy như độ ác tính của khối u, còn $y$ là khuyến nghị điều trị. Bác sĩ nhìn vào
bảng $\hat c_k$ và thấy mô hình đang đọc ra những gì; nếu một giá trị sai thì sửa
nó rồi chạy lại đầu phân loại. Toàn bộ sức hấp dẫn của họ mô hình này nằm trong
hai câu vừa rồi, và phần còn lại của chương là chuyện hai câu ấy đúng tới đâu.

\index{attribution thành phần!khái niệm là thành phần được đặt tên}%
Đặt cạnh khung của chương~\ref{ch:khung-hop-nhat} thì thấy ngay lối thoát này
làm gì. Mục~\ref{sec:ch07-ba-doi-tuong} chia attribution thành ba bài toán theo
đối tượng được chấm điểm, và bài toán thứ ba chấm điểm cho các thành phần $c_k$
bên trong mô hình. Khó khăn của bài toán ấy chưa bao giờ là tính điểm số mà là
đọc nó: một thành phần bên trong một mạng đã huấn luyện không có tên, nên biết
nó quan trọng cũng không nói lên điều gì. Nút thắt khái niệm đảo ngược thứ tự.
Nó không đi tìm tên cho thành phần có sẵn mà dựng thành phần theo tên có sẵn,
nên ký hiệu $c_k$ ở đây là đúng ký hiệu của chương~\ref{ch:khung-hop-nhat} chứ
không phải một chữ trùng nhau: cùng một đối tượng, thêm ràng buộc rằng nó phải
khớp một khái niệm đã được gán nhãn.

\index{mô hình nút thắt khái niệm!ba cách huấn luyện}%
Có ba cách huấn luyện hai chặng ấy, và ba cách cho ba mô hình khác
nhau~\autocite{p30leakage}. Huấn luyện độc lập thì bộ mã hóa học từ $c_k$ và đầu
phân loại học từ $c_k$, mỗi bên một mất mát riêng, không bên nào thấy bên kia.
Huấn luyện tuần tự thì bộ mã hóa học trước như trên, còn đầu phân loại học từ
$\hat c_k$ mà bộ mã hóa đã học xong trả ra. Huấn luyện đồng thời thì cả hai chặng
tối ưu chung một mất mát,
\[
  \beta \, \mathcal{L}_c(\hat c, c) + \mathcal{L}_y(\hat y, y) ,
\]
với $\mathcal{L}_c$ là mất mát trên khái niệm, $\mathcal{L}_y$ là mất mát trên
nhãn, và $\beta \ge 0$ là trọng số quyết định việc học khái niệm nặng bao nhiêu
so với việc học nhiệm vụ. Bài báo viết trọng số ấy là $\lambda$, chữ mà
chương~\ref{ch:giai-thich-duoi-tan-cong} đã cấp cho hệ số chiết khấu nhân quả.
Đặt $\beta = 0$ thì mất mát trên khái niệm biến mất và mô hình trở thành một mô
hình đầu cuối bình thường có một tầng thắt ở giữa.

\index{mô hình nút thắt khái niệm!ba cách mã hóa}%
Chọn cách huấn luyện xong thì còn một lựa chọn nữa, là $\hat c_k$ mang giá trị
gì. Ba lựa chọn thông dụng là bit nhị phân, xác suất trong khoảng từ 0 đến 1, và
logit, tức số thực không chặn. Bài của mục~\ref{sec:ch15-do-ro-ri} đặt tên cho ba
tổ hợp mà nó dùng suốt: mô hình cứng là huấn luyện độc lập cộng mã hóa nhị phân,
mô hình mềm là huấn luyện đồng thời cộng mã hóa xác suất, và mô hình logit là
huấn luyện đồng thời cộng mã hóa logit. Ba tên ấy sẽ trở lại ở mọi mục sau, nên
chúng gộp cả cách huấn luyện lẫn cách mã hóa chứ không chỉ nói
về kiểu số.

\index{mô hình nút thắt khái niệm!hai nhánh có giám sát và không giám sát}%
Cả ba tổ hợp ấy nằm ở một đầu của một dải rộng hơn. Khảo sát của
mục~\ref{sec:ch15-phan-con-lai} chia họ mô hình khái niệm thành hai nhánh theo
việc có nhãn khái niệm hay không~\autocite{p31cbmrisks}. Nhánh có giám sát là
nhánh vừa mô tả: khái niệm do người gán nhãn, nên người ta biết mỗi đơn vị của
tầng thắt phải nhận cái gì, đổi lại là chi phí gán nhãn. Nhánh không giám sát
học các khái niệm thẳng từ nhiệm vụ, không có nhãn khái niệm nào, nên rẻ và mở
rộng được, đổi lại là không ai đặt trước ý nghĩa cho một đơn vị nào và việc diễn
giải chúng để lại cho người dùng. Hai bài đầu của chương làm việc trên nhánh thứ
nhất; mục~\ref{sec:ch15-doc-duoc-dieu-khien-duoc} chuyển sang nhánh thứ hai, và
sẽ thấy chỗ hỏng đổi hình dạng chứ không biến mất.

\index{mô hình nút thắt khái niệm!kiến trúc|)}%
Một chi tiết cuối, vì nó là chỗ mà cả chương xoay quanh. Trong ba cách huấn
luyện, chỉ cách thứ nhất tách hẳn việc học khái niệm khỏi việc học nhiệm vụ. Ở
hai cách kia, mất mát trên nhãn có đường tác động lên bộ mã hóa: đặt giá trị nào
vào $\hat c_k$ để đầu phân loại đoán đúng $y$ là một câu hỏi mà phép tối ưu được
phép trả lời. Mục sau đọc câu trả lời mà nó đưa ra.
